% ==========================================
% CONFIGURACIÓN DE LA PLANTILLA
% ==========================================
%
% ADVERTENCIA: No edites este archivo a menos que
% sepas lo que haces. La mayoría de las
% personalizaciones se hacen desde:
%   - capitulos/portada.tex  (datos del proyecto)
%   - config/proyecto.tex    (fuente, tamaños)
%
% ==========================================

% ==========================================
% PAQUETES ESENCIALES
% ==========================================

% --- Tipografía (XeLaTeX) ---
\usepackage{fontspec}

% Seleccionar fuente según \fuentePrincipal
\makeatletter
\newcommand{\@checkfont}{}
\let\@checkfont\fuentePrincipal
\newcommand{\@Alegreya}{Alegreya}
\newif\ifalegreya
\ifx\@checkfont\@Alegreya
  \alegreyatrue
\else
  \alegreyafalse
\fi
\makeatother

\ifalegreya
  \setmainfont{Alegreya}[
    SmallCapsFont={Alegreya SC},
    Numbers=OldStyle
  ]
  \newfontfamily{\AlegreyaSC}{Alegreya SC}
\else
  \setmainfont{Times New Roman}
  \newcommand{\AlegreyaSC}{\relax}
\fi

% Pasar opciones a xcolor
\PassOptionsToPackage{table,xcdraw}{xcolor}

% --- Idioma ---
\usepackage{csquotes}
\usepackage[spanish,es-tabla]{babel}

% --- Márgenes ---
\usepackage[top=2.5cm, bottom=2.5cm, left=3cm, right=2.5cm]{geometry}

% --- Gráficos e imágenes ---
\usepackage{graphicx}
\usepackage{adjustbox}
\usepackage{float}
\usepackage{caption}
\usepackage{subcaption}

\graphicspath{{imagenes/}}

% --- Colores ---
\usepackage{xcolor}

% --- Tablas ---
\usepackage{booktabs}
\usepackage{longtable}
\usepackage{multirow}

% --- Enlaces y referencias ---
\usepackage[hidelinks]{hyperref}
\usepackage{url}
\usepackage{seqsplit}

% --- Código fuente ---
\usepackage{listings}
\renewcommand{\lstlistingname}{Listado}
\lstset{
  inputencoding=utf8,
  extendedchars=true,
  literate={á}{{\'a}}1 {é}{{\'e}}1 {í}{{\'i}}1 {ó}{{\'o}}1 {ú}{{\'u}}1
  {Á}{{\'A}}1 {É}{{\'E}}1 {Í}{{\'I}}1 {Ó}{{\'O}}1 {Ú}{{\'U}}1
  {ñ}{{\~n}}1 {Ñ}{{\~N}}1 {¿}{{?`}}1 {¡}{{!`}}1,
  breaklines=true,
  breakatwhitespace=true,
  columns=fullflexible,
  keepspaces=true
}

% --- Encabezados y pies ---
\usepackage{fancyhdr}

% --- Espacio entre párrafos ---
\usepackage{parskip}

% --- Formato de títulos ---
\usepackage{titlesec}

\titleformat{\chapter}[display]
  {\normalfont\bfseries}{\chaptertitlename\ \thechapter}{20pt}
  {\fontsize{\tamanoTituloUno}{\dimexpr\tamanoTituloUno*12/10\relax}\selectfont}

\ifalegreya
  \titleformat{\section}
    {\AlegreyaSC\normalfont\bfseries}{\thesection}{1em}
    {\fontsize{\tamanoTituloDos}{\dimexpr\tamanoTituloDos*12/10\relax}\selectfont}
  \titleformat{\subsection}
    {\AlegreyaSC\normalfont\bfseries}{\thesubsection}{1em}
    {\fontsize{\tamanoTituloTres}{\dimexpr\tamanoTituloTres*12/10\relax}\selectfont}
\else
  \titleformat{\section}
    {\normalfont\bfseries}{\thesection}{1em}
    {\fontsize{\tamanoTituloDos}{\dimexpr\tamanoTituloDos*12/10\relax}\selectfont}
  \titleformat{\subsection}
    {\normalfont\bfseries}{\thesubsection}{1em}
    {\fontsize{\tamanoTituloTres}{\dimexpr\tamanoTituloTres*12/10\relax}\selectfont}
\fi

% --- Listas ---
\usepackage{enumitem}

% --- Matemáticas ---
\usepackage{amsmath}
\usepackage{amsfonts}
\usepackage{amssymb}

% --- Apéndices ---
\usepackage[titletoc]{appendix}

% --- Bibliografía ---
\usepackage[backend=biber, style=ieee, sorting=none]{biblatex}
\addbibresource{referencias.bib}

% ==========================================
% CONFIGURACIÓN DE ESTILOS
% ==========================================

% --- Colores para código ---
\definecolor{codegreen}{rgb}{0,0.6,0}
\definecolor{codegray}{rgb}{0.5,0.5,0.5}
\definecolor{codepurple}{rgb}{0.58,0,0.82}
\definecolor{backcolour}{rgb}{0.95,0.95,0.92}

% --- Estilo de listings ---
\lstdefinestyle{mystyle}{
    backgroundcolor=\color{backcolour},
    commentstyle=\color{codegreen},
    keywordstyle=\color{magenta},
    numberstyle=\tiny\color{codegray},
    stringstyle=\color{codepurple},
    basicstyle=\ttfamily\footnotesize,
    breakatwhitespace=false,
    breaklines=true,
    captionpos=b,
    keepspaces=true,
    numbers=left,
    numbersep=5pt,
    showspaces=false,
    showstringspaces=false,
    showtabs=false,
    tabsize=2
}
\lstset{style=mystyle}

% --- Encabezados ---
\setlength{\headheight}{14.2pt}
\pagestyle{fancy}
\fancyhf{}
\fancyhead[L]{\leftmark}
\fancyhead[R]{\thepage}
\renewcommand{\headrulewidth}{0.5pt}

% ==========================================
% COMANDOS ÚTILES
% ==========================================

\newcommand{\comillas}[1]{``#1''}
\newcommand{\codigo}[1]{\colorbox{backcolour}{\texttt{#1}}}
\newcommand{\nota}[1]{\marginpar{\footnotesize\textit{#1}}}
\newcommand{\propia}{Fuente: Elaboración propia.}
